\textbf{Phần 1. Trắc nghiệm lựa chọn}
\setcounterpageref{ex}{0}

\begin{ex} %[1D8V31]
Tính kết quả của phép tính: $47 - 5$.
\choice
{$41$}
{$43$}
{\True $42$}
{$40$}
\loigiai{
Hàng đơn vị: $7 - 5 = 2$. Giữ nguyên 4 chục: $47 - 5 = 42$.
\textbf{Đáp số:} $42$.\par
\textbf{Đáp án của câu hỏi là C}
}
\end{ex}

\begin{ex} %[1D8V31]
Tính kết quả của phép tính: $68 - 6$.
\choice
{\True $62$}
{$61$}
{$63$}
{$60$}
\loigiai{
Hàng đơn vị: $8 - 6 = 2$. Giữ nguyên 6 chục: $68 - 6 = 62$.
\textbf{Đáp số:} $62$.\par
\textbf{Đáp án của câu hỏi là A}
}
\end{ex}

\begin{ex} %[1D8V31]
Tính kết quả của phép tính: $89 - 4$.
\choice
{$84$}
{$86$}
{$83$}
{\True $85$}
\loigiai{
Hàng đơn vị: $9 - 4 = 5$. Giữ nguyên 8 chục: $89 - 4 = 85$.
\textbf{Đáp số:} $85$.\par
\textbf{Đáp án của câu hỏi là D}
}
\end{ex}

\begin{ex} %[1D8V31]
Điền số thích hợp vào ô trống: $56 - \dots = 51$.
\choice
{$4$}
{\True $5$}
{$6$}
{$3$}
\loigiai{
Muốn tìm số trừ, lấy: $56 - 51 = 5$.
\textbf{Đáp số:} $5$.\par
\textbf{Đáp án của câu hỏi là B}
}
\end{ex}

\begin{ex} %[1D8V31]
Điền số thích hợp vào ô trống: $\dots - 3 = 74$.
\choice
{$76$}
{$78$}
{\True $77$}
{$75$}
\loigiai{
Muốn tìm số bị trừ, lấy: $74 + 3 = 77$.
\textbf{Đáp số:} $77$.\par
\textbf{Đáp án của câu hỏi là C}
}
\end{ex}

\begin{ex} %[1D8V31]
Tính giá trị của biểu thức: $78 - 3 - 2$.
\choice
{\True $73$}
{$72$}
{$74$}
{$71$}
\loigiai{
Thực hiện: $78 - 3 = 75$; $75 - 2 = 73$.
\textbf{Đáp số:} $73$.\par
\textbf{Đáp án của câu hỏi là A}
}
\end{ex}

\begin{ex} %[1D8V31]
Tính giá trị của biểu thức: $96 - 4 - 2$.
\choice
{$80$}
{$100$}
{$89$}
{\True $90$}
\loigiai{
Thực hiện: $96 - 4 = 92$; $92 - 2 = 90$.
\textbf{Đáp số:} $90$.\par
\textbf{Đáp án của câu hỏi là D}
}
\end{ex}

\begin{ex} %[1D8V31]
Lớp 1B có 38 học sinh, trong đó có 5 học sinh vắng mặt. Hỏi lớp 1B có bao nhiêu học sinh đang có mặt ở lớp?
\choice
{$32$}
{\True $33$}
{$34$}
{$31$}
\loigiai{
Số học sinh có mặt là: $38 - 5 = 33$ (học sinh).
\textbf{Đáp số:} $33$.\par
\textbf{Đáp án của câu hỏi là B}
}
\end{ex}

\begin{ex} %[1D8V31]
Một cây dừa có 29 quả bắp dừa. Bác nông dân hái bớt 7 quả. Hỏi trên cây dừa còn lại bao nhiêu quả?
\choice
{$21$}
{$23$}
{\True $22$}
{$20$}
\loigiai{
Số quả dừa còn lại là: $29 - 7 = 22$ (quả).
\textbf{Đáp số:} $22$.\par
\textbf{Đáp án của câu hỏi là C}
}
\end{ex}

\begin{ex} %[1D8V31]
Một cửa hàng có 68 chiếc xe máy, đã bán được 5 chiếc xe máy. Hỏi cửa hàng còn lại bao nhiêu chiếc xe máy?
\choice
{\True $63$}
{$62$}
{$64$}
{$61$}
\loigiai{
Số xe máy còn lại là: $68 - 5 = 63$ (chiếc xe máy).
\textbf{Đáp số:} $63$.\par
\textbf{Đáp án của câu hỏi là A}
}
\end{ex}

\begin{ex} %[1D8V31]
Tìm số tự nhiên $x$ thỏa mãn: $87 - x = 82$.
\choice
{$4$}
{$6$}
{$3$}
{\True $5$}
\loigiai{
Ta có: $x = 87 - 82 = 5$.
\textbf{Đáp số:} $5$.\par
\textbf{Đáp án của câu hỏi là D}
}
\end{ex}

\begin{ex} %[1D8V31]
Tìm số tự nhiên $y$ thỏa mãn: $y - 7 = 41$.
\choice
{$47$}
{\True $48$}
{$49$}
{$46$}
\loigiai{
Ta có: $y = 41 + 7 = 48$.
\textbf{Đáp số:} $48$.\par
\textbf{Đáp án của câu hỏi là B}
}
\end{ex}

\textbf{Phần 2. Trắc nghiệm đúng sai}
\setcounterpageref{ex}{0}

\begin{ex} %[1D8V31]
Xét tính đúng hoặc sai của các khẳng định toán học sau đây:
\choiceTF
{\True Một sợi dây dài 59 cm. Bạn Linh cắt bớt một đoạn dài 6 cm. sợi dây còn lại dài $53\text{ cm}$ xăng-ti-mét.}
{Quyển truyện có 76 trang. Bạn Nam đã đọc được 4 trang. Nam còn phải đọc $73$ trang nữa.}
{\True Kết quả của bài toán «Tìm một số biết rằng lấy 88 trừ đi số đó thì được 81» là $7$.}
{Hai số có hiệu bằng 4, số lớn là 69. số bé là số $64$ .}
\loigiai{
\begin{itemize}
    \item Ý a) Đúng: Chiều dài sợi dây còn lại là: $59 - 6 = 53$ (cm).
\textbf{Đáp số:} $53$.
    \item Ý b) Sai: Kết quả đúng là $72$, không phải $73$. Số trang Nam còn phải đọc là: $76 - 4 = 72$ (trang).
\textbf{Đáp số:} $72$.
    \item Ý c) Đúng: Số cần tìm là: $88 - 81 = 7$.
\textbf{Đáp số:} $7$.
    \item Ý d) Sai: Kết quả đúng là $65$, không phải $64$. Số bé là: $69 - 4 = 65$.
\textbf{Đáp số:} $65$.
\end{itemize}
\begin{itemize}
    \item a) \textbf{Đúng}.
    \item b) \textbf{Sai}.
    \item c) \textbf{Đúng}.
    \item d) \textbf{Sai}.
\end{itemize}
}
\end{ex}

\begin{ex} %[1D8V31]
Xét tính đúng hoặc sai của các bài toán và phép tính sau đây:
\choiceTF
{Kết quả của bài toán «Số 58 trừ đi số nào để được 52» là $7$.}
{\True Giá trị của phép tính $49 - 2 - 4$ là $43$.}
{Vườn nhà bà có 48 cây bắp cải, bà đã thu hoạch 6 cây. trong vườn còn lại $43$ cây bắp cải.}
{\True Kết quả của bài toán «Tìm một số biết rằng lấy 88 trừ đi số đó thì được 81» là $7$.}
\loigiai{
\begin{itemize}
    \item Ý a) Sai: Kết quả đúng là $6$, không phải $7$. Số đó là: $58 - 52 = 6$.
\textbf{Đáp số:} $6$.
    \item Ý b) Đúng: Thực hiện: $49 - 2 = 47$; $47 - 4 = 43$.
\textbf{Đáp số:} $43$.
    \item Ý c) Sai: Kết quả đúng là $42$, không phải $43$. Số cây bắp cải còn lại là: $48 - 6 = 42$ (cây).
\textbf{Đáp số:} $42$.
    \item Ý d) Đúng: Số cần tìm là: $88 - 81 = 7$.
\textbf{Đáp số:} $7$.
\end{itemize}
\begin{itemize}
    \item a) \textbf{Sai}.
    \item b) \textbf{Đúng}.
    \item c) \textbf{Sai}.
    \item d) \textbf{Đúng}.
\end{itemize}
}
\end{ex}

\textbf{Phần 3. Trả lời ngắn}
\setcounterpageref{ex}{0}

\begin{ex} %[1D8V31]
Một đàn gà có 39 con gà. Có 7 con gà chạy ra vườn kiếm ăn. Hỏi trong chuồng còn lại bao nhiêu con gà?
\shortans{32}
\loigiai{
\textbf{Đáp số:} $32$\par
Số con gà trong chuồng còn lại là: $39 - 7 = 32$ (con).
\textbf{Đáp số:} $32$.
}
\end{ex}

\begin{ex} %[1D8V31]
Tìm số tự nhiên $a$ sao cho: $a + 4 = 79 - 1$.
\shortans{74}
\loigiai{
\textbf{Đáp số:} $74$\par
Vế phải: $79 - 1 = 78$. Ta có: $a + 4 = 78 \Rightarrow a = 78 - 4 = 74$.
\textbf{Đáp số:} $74$.
}
\end{ex}

\begin{ex} %[1D8V31]
So sánh hai biểu thức: $67 - 5$ và $68 - 4$. Kết quả của biểu thức lớn hơn bằng bao nhiêu?
\shortans{64}
\loigiai{
\textbf{Đáp số:} $64$\par
$67 - 5 = 62$; $68 - 4 = 64$. Giá trị lớn hơn là $64$.
\textbf{Đáp số:} $64$.
}
\end{ex}

\begin{ex} %[1D8V31]
Bạn Hùng có 46 que tính, dùng 4 que xếp hình vuông. Hỏi Hùng còn lại bao nhiêu que tính?
\shortans{42}
\loigiai{
\textbf{Đáp số:} $42$\par
Số que tính Hùng còn lại là: $46 - 4 = 42$ (que tính).
\textbf{Đáp số:} $42$.
}
\end{ex}

\begin{ex} %[1D8V31]
Một bể nước có 89 lít nước. Người ta dùng hết 8 lít nước. Hỏi trong bể còn lại bao nhiêu lít nước?
\shortans{81}
\loigiai{
\textbf{Đáp số:} $81$\par
Số lít nước còn lại là: $89 - 8 = 81$ (lít).
\textbf{Đáp số:} $81$.
}
\end{ex}

\begin{ex} %[1D8V31]
Tính: $98 - 1 - 2 - 3$.
\shortans{92}
\loigiai{
\textbf{Đáp số:} $92$\par
Thực hiện: $98 - 1 = 97$; $97 - 2 = 95$; $95 - 3 = 92$.
\textbf{Đáp số:} $92$.
}
\end{ex}
